% Кодировка и язык
\usepackage[utf8]{inputenc}           % Кодировка исходного файла (UTF-8)
\usepackage[T2A]{fontenc}             % Кодировка шрифтов для кириллицы
\usepackage[english,russian]{babel}   % Локализация: русские названия, переносы

% Поиск в PDF
\usepackage{cmap}                     

% Поля страницы
\usepackage{geometry}
\geometry{left=2cm, right=2cm, top=2cm, bottom=2cm}

% Математика
\usepackage{amsmath, amssymb, amsthm} % Базовые пакеты для математики
\usepackage{mathtools}                % Расширение amsmath (доп. символы и окружения)

% Картинки
\usepackage{graphicx}                 % Вставка картинок (\includegraphics)
\graphicspath{{pictures/}}            % Путь к папке с картинками
\usepackage{float}                    % Позволяет использовать [H] для жёсткого позиционирования
\usepackage{wrapfig}                  % Обтекание текста вокруг картинок

% Таблицы
\usepackage{array}                    % Расширенные возможности для таблиц
\usepackage{booktabs}                 % Красивые таблицы (\toprule, \midrule, \bottomrule)
\usepackage{multirow}                 % Объединение строк в таблицах
\usepackage{tabularx}                 % Таблицы с автоматической шириной колонок
\usepackage{ragged2e}                 % Улучшенное выравнивание (используется в tabularx)

% Списки
\usepackage{enumitem}                 % Гибкая настройка списков (itemize, enumerate)

% Колонтитулы
\usepackage{fancyhdr}                 

% Гиперссылки
\usepackage{xcolor}                   % Работа с цветами
\usepackage{hyperref}                 % Гиперссылки, оглавление, перекрёстные ссылки

% Умные ссылки
\usepackage{cleveref}                 

% Типографика
\usepackage{microtype}                % Улучшает переносы и межсловные расстояния

% Подписи к картинкам и таблицам
\usepackage{caption}                  % Настройка подписей (шрифт, отступы)

% Рисование схем и графиков
\usepackage{tikz}                     % Рисование схем, графиков, диаграмм
\usepackage{pgfplots}                 % Построение графиков функций (на основе tikz)
\pgfplotsset{compat=newest}           % Версия совместимости pgfplots

% Алгоритмы и код
\usepackage{algorithm}                % Окружения для алгоритмов
\usepackage{algorithmic}              % Синтаксис для алгоритмов
\usepackage{listings}                 % Вставка кода с подсветкой синтаксиса


% Настройка нумерации секций и подсекций
\renewcommand{\thesection}{\arabic{section}}                   
\renewcommand{\thesubsection}{\thesection.\arabic{subsection}} 


% Настройка списков: убираем лишние отступы слева
\setlist[itemize]{leftmargin=*}
\setlist[enumerate]{leftmargin=*}


% Определение окружений для теорем, лемм, определений и т.д.
% Нумерация привязана к номеру секции
\newtheorem{theorem}{Теорема}[section]
\newtheorem{definition}{Определение}[section]
\newtheorem{remark}{Замечание}[section]
\newtheorem{lemma}{Лемма}[section]
\newtheorem{example}{Пример}[section]
\newtheorem{consequence}{Следствие}[section]
\newtheorem{algorithmm}{Алгоритм}
\newtheorem{question}{Вопрос}[section]
\newtheorem{answer}{Ответ}[section]
\newtheorem{statement}{Утверждение}[section]
\newtheorem{exercise}{Упражнение}[section]


% Определение цветов
\definecolor{pblue}{rgb}{0.13,0.13,1}    
\definecolor{pgreen}{rgb}{0,0.5,0}       
\definecolor{pred}{rgb}{0.9,0,0}         
\definecolor{pgrey}{rgb}{0.46,0.45,0.48}